% arara: xelatex
\documentclass[12pt]{article}

% \usepackage{physics} % это трэшовый пакет!

\usepackage{hyperref}
\hypersetup{
    colorlinks=true,
    linkcolor=blue,
    filecolor=magenta,      
    urlcolor=cyan,
    pdftitle={Overleaf Example},
    pdfpagemode=FullScreen,
    }
\urlstyle{same} % шрифт из основного текста

\usepackage{verse}

\usepackage{ulem} % перечёркнутый шрифт \sout{что-то} 

\usepackage{tikzducks}

\usepackage{tikz} % картинки в tikz
\usetikzlibrary{arrows.meta}
\usetikzlibrary{shapes, arrows, positioning}
\usepackage{tikz-3dplot} 
\usepackage{tkz-euclide}

\usepackage{verbatim}
\usepackage{microtype} % свешивание пунктуации

\usepackage{array} % для столбцов фиксированной ширины

\usepackage{indentfirst} % отступ в первом параграфе

\usepackage{sectsty} % для центрирования названий частей
\allsectionsfont{\centering}

\usepackage{amsmath, amsfonts, amssymb, amsthm} % куча стандартных математических плюшек

\usepackage{comment}

\usepackage[top=2cm, left=1.2cm, right=1.2cm, bottom=2cm]{geometry} % размер текста на странице

\usepackage{lastpage} % чтобы узнать номер последней страницы

\usepackage{enumitem} % дополнительные плюшки для списков
%  например \begin{enumerate}[resume] позволяет продолжить нумерацию в новом списке
\usepackage{caption}

\usepackage{url} % to use \url{link to web}

\newcommand{\smallduck}{\begin{tikzpicture}[scale=0.3]
    \duck[
        cape=black,
        hat=black,
        mask=black
    ]
    \end{tikzpicture}}

\usepackage{fancyhdr} % весёлые колонтитулы
\pagestyle{fancy}
\lhead{}
\chead{}
\rhead{}
\lfoot{}
\cfoot{}
\rfoot{}

\renewcommand{\headrulewidth}{0.4pt}
\renewcommand{\footrulewidth}{0.4pt}

\usepackage{tcolorbox} % рамочки!
\usepackage{graphicx} % рисунки
\usepackage{todonotes} % для вставки в документ заметок о том, что осталось сделать
% \todo[inline]{Здесь надо коэффициенты исправить}
% \missingfigure{Здесь будет Последний день Помпеи}
% \listoftodos - печатает все поставленные \todo'шки


% более красивые таблицы
\usepackage{booktabs}
% заповеди из докупентации:
% 1. Не используйте вертикальные линни
% 2. Не используйте двойные линии
% 3. Единицы измерения - в шапку таблицы
% 4. Не сокращайте .1 вместо 0.1
% 5. Повторяющееся значение повторяйте, а не говорите "то же"


\setcounter{MaxMatrixCols}{20}
% by crazy default pmatrix supports only 10 cols :)


\usepackage{fontspec}
\usepackage{libertine}
\usepackage{polyglossia}

\setmainlanguage{english}
\setotherlanguages{russian}

% download "Linux Libertine" fonts:
% http://www.linuxlibertine.org/index.php?id=91&L=1
% \setmainfont{Linux Libertine O} % or Helvetica, Arial, Cambria
% why do we need \newfontfamily:
% http://tex.stackexchange.com/questions/91507/
% \newfontfamily{\cyrillicfonttt}{Linux Libertine O}

%\AddEnumerateCounter{\asbuk}{\russian@alph}{щ} % для списков с русскими буквами
%\setlist[enumerate, 1]{label=\asbuk*),ref=\asbuk*}



\usepackage[bibencoding=auto, backend=biber, sorting=none, style=alphabetic]{biblatex}
\addbibresource{cgt.bib}

\graphicspath{{./figures}}

%% эконометрические сокращения
\DeclareMathOperator{\Cov}{\mathbb{C}ov}
\DeclareMathOperator{\Corr}{\mathbb{C}orr}
\DeclareMathOperator{\sCov}{sCov}
\DeclareMathOperator{\sCorr}{sCorr}
\DeclareMathOperator{\sVar}{sVar}
\DeclareMathOperator{\Var}{\mathbb{V}ar}
\DeclareMathOperator{\hVar}{\widehat{\Var}}
\DeclareMathOperator{\hCov}{\widehat{\Cov}}
\DeclareMathOperator{\hCorr}{\widehat{\Corr}}
\DeclareMathOperator{\col}{col}
\DeclareMathOperator{\colspan}{colspan}

\DeclareMathOperator{\se}{se}
\DeclareMathOperator{\row}{row}
\DeclareMathOperator{\grad}{grad}
\DeclareMathOperator{\BestLin}{BestLin}

\let\P\relax
\DeclareMathOperator{\P}{\mathbb{P}}

\DeclareMathOperator{\E}{\mathbb{E}}
\DeclareMathOperator{\trace}{trace}
\DeclareMathOperator{\rank}{rank}
\DeclareMathOperator{\tr}{\trace}
\DeclareMathOperator{\rk}{\rank}
\DeclareMathOperator{\card}{card}
\DeclareMathOperator{\Span}{span}
% \span используется латехом для внутренних нужд

\DeclareMathOperator{\mgf}{mgf}

\DeclareMathOperator{\Convex}{Convex}
\DeclareMathOperator{\plim}{plim}

\usepackage{mathtools}
\DeclarePairedDelimiter{\norm}{\lVert}{\rVert}
\DeclarePairedDelimiter{\abs}{\lvert}{\rvert}
\DeclarePairedDelimiter{\scalp}{\langle}{\rangle}
\DeclarePairedDelimiter{\ceil}{\lceil}{\rceil}

% https://tex.stackexchange.com/questions/12910/how-do-i-typeset-vertical-and-horizontal-lines-inside-a-matrix
% для линий в матрицах
\newcommand*{\horzbar}{\rule[.5ex]{2.5ex}{0.5pt}}
\newcommand*{\vertbar}{\rule[-1ex]{0.5pt}{2.5ex}}


\newcommand{\cN}{\mathcal{N}}
\newcommand{\cF}{\mathcal{F}}

\newcommand{\RR}{\mathbb{R}}
\newcommand{\NN}{\mathbb{N}}

\newcommand{\dBern}{\mathrm{Bern}}
\newcommand{\dPois}{\mathrm{Pois}}
\newcommand{\dBin}{\mathrm{Bin}}
\newcommand{\dMult}{\mathrm{Mult}}
\newcommand{\dGeom}{\mathrm{Geom}}
\newcommand{\dNHGeom}{\mathrm{NHGeom}}
\newcommand{\dHGeom}{\mathrm{HGeom}}
\newcommand{\dDUnif}{\mathrm{DUnif}}
\newcommand{\dFS}{\mathrm{FS}}
\newcommand{\dNBin}{\mathrm{NBin}}

\newcommand{\dTri}{\mathrm{Triangle}}
\newcommand{\dUnif}{\mathrm{Unif}}
\newcommand{\dCauchy}{\mathrm{Cauchy}}
\newcommand{\dN}{\mathcal{N}}
\newcommand{\dLN}{\mathcal{LN}}
\newcommand{\dExpo}{\mathrm{Expo}}
\newcommand{\dBeta}{\mathrm{Beta}}
\newcommand{\dGamma}{\mathrm{Gamma}}
\newcommand{\dWei}{\mathrm{Wei}}
\newcommand{\dLogistic}{\mathrm{Logistic}}
\newcommand{\dRayleigh}{\mathrm{Rayleigh}}
\newcommand{\dPareto}{\mathrm{Pareto}}


% \renewcommand{\b}{\beta} % commented to avoid conflicts with base latex
% \renewcommand{\u}{u} % commented to avoid conflicts with base latex




% выделение питоновских и других пакетов
\newcommand{\sklearn}{sklearn}
\newcommand{\statsmodels}{statsmodels}


\usepackage{thmtools} % тонкая настройка окружения теорем


% окружения!
\declaretheorem[
style=definition, 
%thmbox=M,
title=Theorem, 
numberwithin=section, 
]{theorem}


\declaretheorem[
style=definition, 
title=Lemma, 
numberwithin=section, 
]{lemma}


\declaretheorem[
style=definition, 
title=Idea, 
% numberwithin=section, % ideas are rare!
]{idea} 



% опция sibling=theorem посередине даёт сквозную нумерацию с теоремами
\declaretheorem[
style=definition, 
title=Definition, 
sibling=theorem, 
]{definition}

\declaretheorem[
style=definition, 
title=Problem, 
% sibling=theorem, 
]{problem}

\declaretheorem[
style=definition, 
title=Exercise, 
% sibling=theorem, 
numberwithin=section,
]{exercise}


\declaretheorem[
style=definition,
title=Solution,
numbered=no
]{sol}


%\newenvironment{problem}{}{} % environment just prints the text
% \newenvironment{sol}{}{}
\newcommand{\trru}[1]{#1}
\newcommand{\tren}[1]{#1}



\begin{document}

% задачи можно дёргать:
% https://github.com/bdemeshev/metrics_pro/blob/master/metrics_body.tex
% и ещё (слегка модифицируя)
% https://github.com/bdemeshev/em_pset
% картиночки и доказательства на английском:
% https://github.com/olyagnilova/gauss-markov-pythagoras

\tableofcontents{}



\section{Classic binomial and multinomial identities}


\begin{definition}
    Factorial of $n$, denoted by $n!$, is the number of ways to put $n$ distinct objects in a queue. 
\end{definition}


\begin{definition}
    Binomial coefficient $n \choose k$ is the number of ways to select $k$ distinct objects from $n$ distinct objects.
\end{definition}

\begin{definition}
    For $n = a + b + c$ the multinomial coefficient $n \choose a, b, c$ is the number of ways to split $n$ distinct objects into different boxes, 
    where the first box contains $a$ objects, the second box contains $b$ objects, the third box contains $c$ objects. 

    One may create multinomial coefficient for any number of boxes :)
\end{definition}


\begin{exercise}
Is $n \choose k$ equal to $n \choose n-k$? 

Is $n \choose k$ equal to $n \choose k, n-k$?
\end{exercise}

\begin{exercise}
Write the product
    \[
    {100 \choose 20} \cdot {80 \choose 50} \cdot {30 \choose 10}
    \]
as one multinomial coefficient.
\end{exercise}

\begin{exercise}
Write 
\[
n \choose a, b, c, d
\] using binomial coefficients.

Write 
\[
n \choose a, b, c, d
\]
using factorials of $n$, $a$, $b$, $c$ and $d$.
\end{exercise}

\begin{idea}
Solution strategies for combinatorial identities:
\begin{enumerate}
    \item Double counting or story proof. 
    \item Visual identity (rare gem!).
    \item Evaluation of a polynomial at a special point. Differentiation, integration.
    \item Combine old identities into new. 
    \item Guess the answer and verify by induction. 
    \item Dark art!
\end{enumerate}
\end{idea}

\subsection*{Pascal's triangle and rule}
\begin{exercise}
Recall the Pascal arithmetic triangle. Draw two versions: triangular and rectangular. 
\end{exercise}

% Here the answer usually contains just one binomial coefficient or no at all:
\begin{exercise}
Pascal's rule. Simplify
\[
{n \choose k} + {n \choose k+1}.
\]
\[
{n \choose a-1, b, c} + {n \choose a, b-1, c} + {n \choose a, b, c-1}.
\]
\[
{n \choose a-1, b, c, d} + {n \choose a, b-1, c, d} + {n \choose a, b, c-1, d} + {n \choose a, b, c, d-1}.
\]
\end{exercise}

\begin{exercise}
Simplify 
\[
{n \choose 0} + {n \choose 1} + {n \choose 2} + \dots + {n \choose n}.
\]
Recall the Pascal's rule and simplify 
\[
{2n \choose 0} + {2n \choose 2} + {2n \choose 4} + \dots + {2n \choose 2n}.
\]
Simplify
\[
{2n-1 \choose 0} + {2n-1 \choose 2} + {2n-1 \choose 4} + \dots + {2n-1 \choose 2n-2}.
\]
\end{exercise}

\begin{exercise}
Simplify 
\[
\sum_{a + b + c = 2026} {2026 \choose a, b, c}.
\]    
\[
\sum_{a + b + c + d = 2027} {2027 \choose a, b, c, d}.
\]    
\[
\sum_{a + b + c = 100} {100 \choose a} \cdot {100-a \choose b}.
\]
\[
\sum_{a + b \leq n} {n \choose a} \cdot {n-a \choose b}.
\]
Here $a$, $b$ and $c$ are assumed non-negative integers.
\end{exercise}

\subsection*{Hockey-stick identity}

\begin{exercise}
Hockey-stick identity, Christmas stocking identity, boomerang identity.

Simplify
\[
\sum_{n = a}^b {n \choose a}.
\]
\[
1 + 2 + 3 + \dots + n.
\]
\[
1\cdot 2 + 2 \cdot 3 + 3 \cdot 4 + \dots + n \cdot (n+1).
\]
\[
1\cdot 2 \cdot 3 + 2 \cdot 3 \cdot 4 + 3 \cdot 4 \cdot 5 + \dots + n \cdot (n+1) \cdot (n+2).
\]
\end{exercise}

\begin{exercise}
Obtain closed form expressions for:     
\[
1^2 + 2^2 + \dots + n^2,
\]
\[
1^3 + 2^3 + \dots + n^3,
\]
\[
1^4 + 2^4 + \dots + n^4.
\]
\end{exercise}

\begin{exercise}
Hockey-sticks plus idea from linear algebra!

Decompose $n^2 + 2n + 5$ using $n \choose 0$, $n \choose 1$, $n \choose 2$ as the basis. 

Find the sum 
\[
\sum_{k=2}^{2026} {k^2 + 2k + 5}.
\]

Using similar trick find the sums
\[
\sum_{k=1}^{n} {5k^2 + 4k}, \quad \sum_{k=1}^{n} {k^3 + 2k^2}.
\]
\end{exercise}

\begin{exercise}
Let's start from some number $n \choose k$ inside the Pascal's triangle. 

In one step we can go either in north-east or north-west direction. 

Find the sum of all numbers we can reach.
\end{exercise}


\subsection*{Wandermonde's identity}

\begin{exercise}
Simplify 
\[
{a \choose 0}{b \choose m} + {a \choose 1}{b \choose m-1} + {a \choose 2}{b \choose m-2} + \dots + {a \choose m}{b \choose 0}.
\]
\[
{a \choose 0}{m \choose 0} + {a \choose 1}{m \choose 1} + {a \choose 2}{m \choose 3} + \dots + {a \choose m}{m \choose m}.
\]
\[
{a \choose 0}^2 + {a \choose 1}^2  + {a \choose 2}^2 + \dots + {a \choose a}^2.
\]
\end{exercise}

\begin{exercise}
There are $n = 500$ fishes in the pond, $t = 30$ of them are tagged. You capture $c = 5$ fishes at random without replacement. 
Let the random variable $Y$ be the number of captured tagged fishes. 
Find the probabilities $\P(Y = 0)$, $\P(Y = 1)$, $\P(Y = 2)$, $\P(Y = y)$.
\end{exercise}

\begin{exercise}
There are $n = 24$ candies in the jar, $14$ with orange flavour and $10$ with ananas flavour.
Your take $6$ candies at random. Let $Y$ be the number of candies with orange flavour you picked. 
Find the probabilities $\P(Y = 0)$, $\P(Y = 1)$, $\P(Y = 2)$, $\P(Y = y)$.
\end{exercise}

\begin{exercise}
A bridge hand (random $13$ cards out of $52$-card standard deck) is dealt. 

What is the probability that the hand contains no spades? 

What is the probability that the hand contains exactly $5$ hearts? 

What is the probability that the hand contains at most two aces? 
\end{exercise}

\begin{exercise}
Simplify 
\[
\sum_{a + b + c = k}{n_1 \choose a}{n_2 \choose b}{n_3 \choose c}. % {n_1 + n_2 + n_3 \choose k}
\]
\end{exercise}

\subsection*{Binomial or Taylor or MacLaurin expansion}
\begin{exercise}
Write the Taylor expansion (binomial expansion) at $x = 0$ for 
\[
(1 + x)^{n}, \quad (x + y)^n.
\]
\end{exercise}


\begin{exercise}
By plugging some special points to the $(x + y)^n$ or somewhere else find the sums
\[
{n \choose 0} + {n \choose 1} + {n\choose 2} + \dots + {n \choose n}.
\]
\[
\sum_{k=0}^n k {n \choose k}, \quad \sum_{k=0}^n (k-1) {n \choose k}, \quad \sum_{k=0}^n (-1)^k {n \choose k}, \quad \sum_{k=0}^n (-1)^{k-1} k {n \choose k}
\]

Hint: differentiation may help you.
\end{exercise}


\subsection*{How that can be?}

\begin{exercise}
Be brave, really brave! 

Take a deep breath and calculate 
\[
{-1 \choose 3}, \quad {-5 \choose 5},\quad  {1.5 \choose 3}, \quad  
{-1/2 \choose 4}, \quad  {1/2 \choose k}, \quad  {-1/2 \choose n}. 
\]
\end{exercise}

\begin{exercise}
Simplify 
\[
{-1/2 \choose 3} + {-1/2 \choose 4}.
\]

Draw the Pascal's triangle starting from $-1/2 \choose 0$. 

Draw the Pascal's triangle starting from $1/3 \choose 0$. 
\end{exercise}


\begin{exercise}
Write the Taylor expansion at $x = 0$ for 
\[
\frac{1}{1 - x}, \quad  (1 + x)^{-5}, \quad  (1 + x)^{-1/2}, \quad  
(1 + 3x)^{1/3}, \quad  (1 - 5x)^{1/2}.
\]
\end{exercise}

\begin{exercise}
Simplify 
\[
{10 \choose 0}{-0.5 \choose 7} + {10 \choose 1}{-0.5 \choose 6} + {10 \choose 2}{-0.5 \choose 5} + \dots + {10 \choose 7}{-0.5 \choose 0}.
\]

Hint: for a honest proof consider the binomial expansion of $(1+x)^10$, $(1+x)^{-0.5}$, $(1+x)^{10-0.5}$.
\end{exercise}


\begin{exercise}
Homage to Newton. 

Write the Taylor expansion at $x=0$ for $(1 - x^2)^{1/2}$.

Visually (or analytically if you wish) find the integral
\[
\int_{-1}^1 (1 - x^2)^{1/2}.
\]

Calculate the sums
\[
{1/2 \choose 0} - \frac{1}{3}{1/2 \choose 1} + \frac{1}{5} {1/2 \choose 2} + \dots
\]
\[
\sum_{k=2}^{\infty} \frac{1}{2k+1} \frac{(2k - 3)!!}{(2k)!!}
\]
\end{exercise}


\subsection*{Gauss summation trick}

\begin{exercise}
Simplify 
\[
1 + 2 + 3 + \dots + n;
\]
\[
1^2 + 2^2 + 3^2 + \dots + n^2;
\]
\[
{n \choose 1} + 2 {n \choose 2} + 3 {n \choose 3} + \dots + n{n \choose n};
\]
\[
{n \choose 0} + 2 {n \choose 1} + 3 {n \choose 2} + \dots + (n+1){n \choose n};
\]
\[
{n \choose 2} + 2 {n \choose 3} + 3 {n \choose 4} + \dots + (n-1){n \choose n};
\]
\[
{n \choose 0} + 3 {n \choose 1} + 5 {n \choose 2} + \dots + (2n+1){n \choose n};
\]
\[
{n \choose 0} - 2 {n \choose 1} + 3 {n \choose 2} - 4 {n \choose 3} \dots + (-1)^n (n+1){n \choose n};
\]
\[
3{n \choose 1} + 7 {n \choose 2} + 11 {n \choose 3} + \dots + (4n-1){n \choose n};
\]
\[
{n \choose 1} - 2 {n \choose 2} + 3 {n \choose 3} - 4 {n \choose 4} + \dots + (-1)^n n{n \choose n};
\]
\end{exercise}


\newpage
First day quiz :)

\begin{enumerate}
\item How many ways are to order $6$ persons in a queue?
\item The box contains $50$ balls numbered from $1$ to $50$. We draw $6$ balls without replacement and write down the resulting sequence of numbers.
How many different sequences are possible?
\item I have $5$ yellow books and $7$ red books. I want to place them on my bookshelf and I would like to place yellow books to the left of red books. 
How many book orders are possible?
\end{enumerate}




\nocite{graham1994concrete}
\nocite{harris2010combinatorics}
\nocite{guichard2025introcgt}    


\newpage
\addcontentsline{toc}{section}{Sources of wisdom}
\printbibliography[heading=none]


\end{document}
